\chapter{Basic Arduino Experiments}

This chapter presents the fundamental Arduino-based experiments carried out during the internship. These experiments introduce basic digital interfacing, display modules, and embedded programming concepts using the Arduino UNO development board. The knowledge gained from these experiments forms the foundation for more advanced embedded system applications discussed in later chapters.

%%%%%%%%%%%%%%%%%%%%%%%%%%%%%%%%%%%%%%%%%%%%%%%%%%%%%%%%%%%%%

\section{LED Flash System}

\subsection{Objective}

The objective of this experiment is to interface an LED with the Arduino UNO and make it blink continuously at a fixed time interval using digital output pins.

\subsection{Problem Statement}

Design and implement a simple LED flashing system using Arduino UNO to understand digital output control, GPIO programming, and timing functions.

\subsection{Components Used}

% Components Table

\subsection{Hardware Connections}

% Pin Connection Table

\subsection{Circuit Diagram}

% Figure Here

\subsection{Software Used}

\begin{itemize}
    \item Arduino IDE
    \item SimulIDE
\end{itemize}

\subsection{Source Code}

The complete Arduino source code for this experiment is available in the project repository.

\href{https://github.com/YOUR_USERNAME/YOUR_REPOSITORY}{\textbf{Click here to view the Arduino source code}}

\subsection{Working Principle}

The Arduino UNO configures digital pin D12 as an output during initialization. The program repeatedly changes the output state of the pin from HIGH to LOW with a one-second delay between each transition. As a result, the connected LED turns ON and OFF continuously, demonstrating the basic operation of digital output control and timing functions in Arduino programming.

\subsection{Results}

The LED blinked successfully with a one-second ON period followed by a one-second OFF period. The experiment verified the correct operation of Arduino digital output pins and demonstrated the use of software-based timing control.

\subsection{Conclusion}

This experiment successfully demonstrated the interfacing of an LED with Arduino UNO. It provided practical understanding of GPIO configuration, digital output control, and program execution flow, forming the basis for subsequent embedded system experiments.